\section{Theme Architecture}
\label{sec:themes}

A theme is \emph{one coherent design system}, defined once at a general, component-agnostic
level. It is \emph{not} an attribute owned separately by each diagram type; it is an
authority that every component conforms to. Each component---flow, sequence, timeline, tree,
class diagram, ER diagram, and so on---provides a \emph{specific implementation} of that
general system: it knows how to realise the theme's tokens in its own geometric terms. The
theme knows nothing about flowcharts or Gantt bars; the component knows nothing about brand
palettes or spacing scales. The theme sets the contract; the component fulfils it.

\begin{quote}
\textbf{Confirmed design principle (binding):} A theme must be a coherent system applicable
to many (ideally all) components. You cannot conceive of a theme as an attribute of each
component as a separate system. A theme is defined once in general terms; each component
provides its own specific implementation of that general definition.
\end{quote}

Themes also drive \textbf{geometry, layout, and routing}---not just colour. A theme defines
widths, fonts, weights, spacing rhythms, density levels, and shape language. The layout
engine consults the theme to make rendering, layout, and routing decisions. The old rule that
``a theme may not alter geometry'' is explicitly retired: it contradicts both this model and
the existing code, where the timeline's \texttt{spineSpacing} and \texttt{nodeWrap} tokens
already drive layout behaviour.

\paragraph{Implementation status.}
As of 2026-06-14 the compiler ships 16 grammar-specific \texttt{theme.ts} files (one per
component: \texttt{FlowTheme}, \texttt{SequenceTheme}, \texttt{ClassTheme},
\texttt{TreeTheme}, etc.) plus the timeline's \texttt{ResolvedTheme} (which covers 14 named
timeline themes). These are separate, incompatible vocabularies---the fragmentation this
section resolves. The migration path is to generalise \texttt{ResolvedTheme} upward into a
shared general contract, then adapt each per-component file to consume that contract.

% ──────────────────────────────────────────────────────────────────────────────
\subsection{Current Reality: Per-Component Fragmentation}
\label{sec:themes-fragmentation}

Today every grammar family carries an independent styling vocabulary:

\begin{itemize}
  \item \textbf{Timeline:} \texttt{packages/core/src/themes/types.ts} ---
        \texttt{ResolvedTheme} with \texttt{canvas}, \texttt{typography},
        \texttt{axis}, \texttt{track}, \texttt{milestone}, \texttt{serpentine},
        \texttt{roadmap}, and layout-affecting flags such as
        \texttt{spineSpacing} and \texttt{nodeWrap}. Fourteen named themes:
        \texttt{consulting}, \texttt{executive}, \texttt{product}, \texttt{release},
        \texttt{minimal}, \texttt{showcase}, \texttt{bytebytego}, \texttt{ai-timeline},
        and others.
  \item \textbf{Flow:} \texttt{FlowTheme}---orientation, edge routing style, node
        geometry, font sizes.
  \item \textbf{Sequence:} \texttt{SequenceTheme}---participant render mode (box vs.\
        card), row heights, activation bar widths, fragment styling.
  \item \textbf{Sixteen further grammars} (class, ER, state, tree, C4, chart,
        gitGraph, journey, sankey, kanban, requirement, block, packet, architecture)
        each with their own isolated token structs.
\end{itemize}

No named theme (``consulting'', ``executive'') spans more than one grammar today. There is no
shared vocabulary for palette roles, type scale, or spacing rhythm across components.
This is the problem the new model resolves.

% ──────────────────────────────────────────────────────────────────────────────
\subsection{The Themes $\times$ Components Matrix}
\label{sec:themes-matrix}

The new model is a two-dimensional system:

\begin{description}
  \item[Theme (the row)] A general design language. It defines palette by role, type
        scale, spacing rhythm, density, and shape language---nothing component-specific.
        Adding a new theme adds a new row.
  \item[Component (the column)] A diagram grammar. It knows how to realise the general
        theme in its own geometric terms. Adding a new component adds a new column.
  \item[Intersection ($T \times C$)] How theme $T$ renders component $C$. This is
        computed by $C$'s implementation consuming $T$'s general tokens; it is not
        stored as a separate artefact.
\end{description}

\begin{table}[h]
\centering
\small
\begin{tabularx}{\textwidth}{l|XXXXXX}
\toprule
\textbf{Theme $\downarrow$ / Component $\rightarrow$} &
  \textbf{Timeline} & \textbf{Flow} & \textbf{Sequence} &
  \textbf{Tree} & \textbf{Class} & \textbf{\ldots} \\
\midrule
Consulting  & \checkmark & \checkmark & \checkmark & \checkmark & \checkmark & \ldots \\
Executive   & \checkmark & \checkmark & \checkmark & \checkmark & \checkmark & \ldots \\
ByteByteGo  & \checkmark & \checkmark & \checkmark & \checkmark & \checkmark & \ldots \\
Minimal     & \checkmark & \checkmark & \checkmark & \checkmark & \checkmark & \ldots \\
\textit{New theme}  & \checkmark & \checkmark & \checkmark & \checkmark & \checkmark & \ldots \\
\bottomrule
\end{tabularx}
\caption{Themes $\times$ Components matrix. A theme is the row authority; a component is the
column implementer. A \checkmark\ means the component's implementation consumes that theme's
tokens. Add a theme once---every component renders it. Add a component once (implement the
contract)---it inherits every existing theme for free.}
\label{tab:themes-matrix}
\end{table}

This is the inverse of the fragmented model. In the fragmented model each component owns its
theme, so adding a new theme requires editing every component; adding a new component
requires writing a new, isolated theme. In the matrix model each direction is O(1) in
authoring work: a new theme is one new general-contract file; a new component is one new
implementation that automatically supports all existing themes.

% ──────────────────────────────────────────────────────────────────────────────
\subsection{The General Theme Contract}
\label{sec:themes-contract}

The general theme contract is the primary design artefact of the theming system. It is an
abstract, component-agnostic token vocabulary expressive enough to drive any component's
geometry, routing, and look---with zero component-specific fields. Every component's
implementation consumes this vocabulary; nothing outside it may be required for a correct
rendering~\cite{dtcg2024}.

\begin{quote}
\textbf{General-contract rule (binding):} The general theme vocabulary MUST be rich enough
that a component's layout engine can ask the theme for spacing, density, routing style, and
palette roles and receive a deterministic answer. No component-specific token may leak into
the general contract. Components may \emph{interpret} tokens in their own context; they may
not \emph{demand} tokens not in the contract.
\end{quote}

The vocabulary is organized into six domains; the specific field names are subject to
refinement but the domain coverage is fixed:

\subsubsection{Palette}
\label{sec:contract-palette}

The general contract carries \textbf{two distinct palette systems}. Structural
components consume the \emph{role palette}; data and chart components consume the
\emph{data palette}. Both systems are part of the general contract.

\paragraph{Role palette (semantic/structural).}
Structural roles are expressed by name, not raw hex. Raw hex values are an
implementation detail of a specific theme instance.

\begin{lstlisting}[language=yaml,caption={General theme contract: role palette (decided)}]
palette:
  # Surface roles
  surface:         "#FFFFFF"   # primary background
  surface-raised:  "#F8F8F8"  # cards, panels
  surface-overlay: "#0A0A0A"  # dark overlay / inverse background
  surface-panel:   "#F4F4F4"  # lane / panel / track-header background (distinct from surface)

  # Content roles
  ink:             "#111111"   # primary text
  ink-muted:       "#666666"   # secondary / de-emphasized text
  ink-inverse:     "#FFFFFF"   # text on dark backgrounds
  ink-panel:       "#333333"   # text/ink rendered on panel or lane backgrounds

  # Brand / accent
  accent:          "#1F497D"   # primary brand colour
  accent-muted:    "#6A92BF"   # lighter accent (inactive, planned)
  accent-strong:   "#0D2B4E"   # darker accent (active, in-progress)

  # Border
  border:          "#DDDDDD"   # default stroke / divider
  border-strong:   "#999999"   # emphasized boundary

  # Muted fill (background regions, disabled states)
  muted:           "#F0F0F0"
  muted-strong:    "#CCCCCC"

  # Semantic status roles (components map these to IR status values)
  status-success:  { fill: "#4CAF50", stroke: "#388E3C" }
  status-warning:  { fill: "#D97706", stroke: "#A35A04" }
  status-error:    { fill: "#DC2626", stroke: "#A81919" }
  status-info:     { fill: "#1F497D", stroke: "#0D2B4E" }

  # IR workflow states (super-set of status roles above)
  status-planned:    { fill: "#6A92BF", stroke: "#1F497D" }
  status-active:     { fill: "#1F497D", stroke: "#0D2B4E" }
  status-done:       { fill: "#888888", stroke: "#555555", opacity: 0.85 }
  status-cancelled:  { fill: "#CCCCCC", stroke: "#999999", opacity: 0.60 }
  status-uncertain:  { fill: "#CCCCCC", stroke: "#999999", opacity: 0.70 }

  # Optional fill pattern applied to workflow-state renders
  status-pattern: "solid"  # solid | hatch | dashed-border
\end{lstlisting}

\paragraph{Data palette (quantities and sequences).}
Chart and data components (pie, XY chart, kanban, gitGraph, radar, sankey) consume the
data palette. It provides three sub-systems covering the full range of quantitative
encodings~\cite{dtcg2024}.

\begin{lstlisting}[language=yaml,caption={General theme contract: data palette (decided)}]
data-palette:
  # Categorical sequence — ordered colours for pie slices, chart series,
  # kanban columns, branch colours. Use by index (0-based); themes may
  # supply any length >= 6.
  categorical:
    - "#1F497D"   # 0 — primary
    - "#6A92BF"   # 1
    - "#4CAF50"   # 2
    - "#D97706"   # 3
    - "#9C27B0"   # 4
    - "#00BCD4"   # 5

  # Sequential ramp — for ordered quantitative encodings (sankey
  # throughput, heat intensity, radar fill). Low → high.
  sequential:
    low:  "#EEF4FB"
    mid:  "#6A92BF"
    high: "#0D2B4E"

  # Diverging ramp — for values around a meaningful midpoint
  # (positive/negative, above/below target). Negative → zero → positive.
  diverging:
    negative: "#DC2626"
    zero:     "#F5F5F5"
    positive: "#4CAF50"
\end{lstlisting}

\subsubsection{Typography}
\label{sec:contract-typography}

Typography specifies family and fallback, a \emph{type scale} (not a list of fixed sizes),
and a weight set. Components consume the scale by step name, not by pixel value.

\begin{lstlisting}[language=yaml,caption={General theme contract: typography block (proposed shape)}]
typography:
  family:    "Helvetica Neue"
  fallback:  "Arial, Helvetica, sans-serif"
  # Embedded font files for deterministic metric computation.
  files:
    regular: "fonts/HelveticaNeue-Regular.woff2"
    bold:    "fonts/HelveticaNeue-Bold.woff2"

  # Type scale: step names → point sizes.
  # Components reference steps by name, not by pt value.
  scale:
    xs:   8    # pt — axis ticks, dense labels
    sm:   9    # pt — captions, secondary labels
    base: 11   # pt — body text, activity labels
    md:   13   # pt — subtitles, section headers
    lg:   18   # pt — document title
    xl:   24   # pt — hero / keynote headline

  # Weight palette
  weights:
    regular: 400
    medium:  500
    semibold: 600
    bold:    700
\end{lstlisting}

\subsubsection{Spacing / Rhythm Scale}
\label{sec:contract-spacing}

A spacing scale defines the rhythm unit and derived named steps. Components consult the
scale as \emph{advice}, not as a binding pixel-identical contract: each component maps
step names onto its own geometry as it sees fit, preserving layout judgement while sharing
a common rhythm. The scale is not normalised across components.

\begin{lstlisting}[language=yaml,caption={General theme contract: spacing block (advisory)}]
spacing:
  unit: 8           # px — base rhythm unit
  # Named steps (advisory; components map these to their own geometry)
  steps:
    xxs:  2   # 0.25u
    xs:   4   # 0.5u
    sm:   8   # 1u
    md:  16   # 2u
    lg:  24   # 3u
    xl:  32   # 4u
    xxl: 48   # 6u
\end{lstlisting}

\subsubsection{Density}
\label{sec:contract-density}

Density is a discrete named level that components use to select between compact and
comfortable configurations. Three levels are defined---no more, no less. The layout engine
consults density to choose row heights, inter-element gaps, and label truncation thresholds
(e.g.\ whether secondary labels are shown at all).

\begin{lstlisting}[language=yaml,caption={General theme contract: density (decided)}]
density: "normal"   # compact | normal | comfortable
\end{lstlisting}

A \texttt{compact} theme instructs every component to select its smallest safe row height,
tightest inter-element gap, and most aggressive label truncation. A \texttt{comfortable}
theme selects the most generous variants. Components map the three density levels to their
own geometry constants; density is not a continuous scalar.

\subsubsection{Shape Language}
\label{sec:contract-shape}

Shape language defines how geometric objects feel. Components use these tokens when
deciding corner radius, node padding, stroke scale, and connector style.

\begin{lstlisting}[language=yaml,caption={General theme contract: shape block (proposed shape)}]
shape:
  corner-radius:  0       # px — 0 = square; >0 = rounded
  node-padding:   8       # px — internal padding inside node shapes
  stroke-scale:   1.0     # multiplier on all stroke widths (0.5 = hairline; 2.0 = bold)
  connector-style: "elbow" # straight | elbow | curved | orthogonal
  marker-shape: "circle"  # circle | diamond | square — milestone / marker glyph
\end{lstlisting}

\subsubsection{Effects and Motion}
\label{sec:contract-effects}

Effects declare the fidelity tier and optional art-effect tokens. These gate which rendering
backend capability profile is required (see Section~\ref{sec:fidelity-tiers}).

\begin{lstlisting}[language=yaml,caption={General theme contract: effects block (proposed shape)}]
effects:
  fidelity: 1        # 0=Minimal | 1=Crisp | 2=Polished | 3=Showcase
  drop-shadow:  { enabled: false }
  glow:         { enabled: false }
  motion:       { enabled: false }   # SMIL animation opt-in (Section~14)
\end{lstlisting}

\subsubsection{The Timeline ResolvedTheme as the Generalisation Seed}
\label{sec:contract-seed}

The timeline's \texttt{ResolvedTheme}
(\texttt{packages/core/src/themes/types.ts}) is the closest existing artefact to the
general contract. It already covers canvas, typography (with weights), axis geometry, track
geometry, and layout-affecting flags
(\texttt{spineSpacing}, \texttt{nodeWrap}, \texttt{maxSubLanes}). It demonstrates that a
theme can---and must---drive layout decisions, not merely colour.

The migration plan is to \emph{generalise upward}: extract the conceptual domains (palette
roles, type scale, spacing rhythm, density, shape language) from \texttt{ResolvedTheme}
into the shared general contract, then recast the timeline-specific fields (axis height,
track header width, serpentine row spacing) as the timeline component's interpretation of
those general tokens~\cite{vegalite2017}.

% ──────────────────────────────────────────────────────────────────────────────
\subsection{Theme Drives Geometry, Layout, and Routing}
\label{sec:themes-drives-layout}

Themes are not a post-processing colour filter. The layout engine queries the theme
\emph{during} layout computation. Representative examples from the existing codebase:

\begin{itemize}
  \item \textbf{\texttt{spineSpacing: 'time'} vs.\ \texttt{'even'}} --- controls whether
        vertical-spine milestone nodes are placed proportionally to elapsed time or at
        uniform y-intervals. This changes node $y$-coordinates, not just colour.
  \item \textbf{\texttt{nodeWrap: 'over-under'}} --- routes the horizontal timeline spine
        around each on-axis node as a semicircular arc (alternating above / below),
        producing entirely different path geometry compared to the default straight line.
  \item \textbf{\texttt{density}} --- compact density reduces \texttt{row\_height} and
        \texttt{sub\_lane\_height}, shrinking the canvas and shifting all track and
        milestone $y$-coordinates.
  \item \textbf{\texttt{shape.corner-radius}} --- a flow diagram component may use this
        token to choose between rectangular and rounded node shapes, which affects label
        area computation and thus label wrap point.
  \item \textbf{\texttt{connector-style: 'orthogonal'}} --- instructs the flow / sequence
        layout engine to use orthogonal edge routing rather than diagonal or curved,
        changing every edge path.
\end{itemize}

\begin{quote}
\textbf{Layout-engine rule (binding):} A layout engine MUST consult the active theme for
any geometry decision that varies across themes. It MUST NOT hard-code values that the
general contract expresses. Theme-driven geometry is still \emph{closed-form}: the theme
supplies constants; the engine computes coordinates deterministically from those constants.
No iterative solver or random sampling is introduced.
\end{quote}

% ──────────────────────────────────────────────────────────────────────────────
\subsection{Reach: Theme, Layout, and New Components}
\label{sec:themes-reach}

Drawing the right boundary between theme, layout argument, and new grammar is the
architectural decision that keeps the system from proliferating components unnecessarily.
The rule is stated at the IR boundary:

\begin{description}
  \item[\emph{New component (grammar)}] Only when the diagram requires \emph{new semantic
        data} that the existing IR cannot carry. A new grammar is a last resort.
  \item[\emph{Layout argument}] When the same data can be positioned differently without
        needing new IR fields. A layout mode is a named argument passed at render time,
        not a new grammar.
  \item[\emph{Theme tokens}] Look-and-feel within a layout: colour, typography, spacing,
        shape, effects. A theme produces a visually distinct output from the same IR +
        same layout.
\end{description}

\paragraph{Worked example: the timeline as six diagrams.}

The timeline IR is one grammar. A single set of entries---each with a date, label, status,
and optional icon---can be rendered six ways that look and behave radically differently,
purely via \texttt{\{layout: $L$, theme: $T$\}}:

\begin{table}[h]
\centering
\small
\begin{tabularx}{\textwidth}{llX}
\toprule
\textbf{Layout argument} & \textbf{Theme} & \textbf{Visual result} \\
\midrule
\texttt{vertical-spine}  & consulting   & Spine-and-card with navy accent, tidy whitespace \\
\texttt{horizontal-arc}  & executive    & Arc path with serif labels, subtle quarter shading \\
\texttt{roadmap}         & product      & Horizontal Gantt-style bars, dense information \\
\texttt{serpentine}      & showcase     & Gradient boustrophedon path, premium glow effects \\
\texttt{gantt}           & release      & Traffic-light Gantt, monospace labels, weekly grid \\
\texttt{timeline-columns}& minimal      & Column-per-period grid, greyscale, LaTeX-safe \\
\bottomrule
\end{tabularx}
\caption{One timeline IR; six visual outputs via \{layout, theme\} arguments. No new grammar
is required for any of these. Same data, radically different look and behaviour.}
\label{tab:reach-example}
\end{table}

The rule of thumb: if the data the IR already carries is sufficient to drive the new visual,
use a layout argument and/or theme tokens. Only introduce a new grammar when the IR would
need to grow new fields to support it (e.g., swimlane headers, dependency arrows, resource
calendars---these are genuinely new semantic data, not layout variations).

% ──────────────────────────────────────────────────────────────────────────────
\subsection{Mermaid Configuration Surface}
\label{sec:themes-mermaid-surface}

All user-facing surface is Mermaid or Mermaid-like. Theme selection, token overrides,
layout selection, and density are expressed via Mermaid-native configuration mechanisms,
so a user can remain in standard Mermaid syntax while accessing the extended vocabulary.

\paragraph{Frontmatter (\texttt{---}~\ldots~\texttt{---}).}

\begin{lstlisting}[language=yaml,caption={Theme and layout selection via Mermaid YAML frontmatter}]
---
theme: executive          # named theme from the registry
layout: vertical-spine    # layout family for this diagram
density: compact          # compact | normal | comfortable
tokens:                   # inline token overrides (scoped to this diagram)
  palette.accent: "#8B0000"
  typography.scale.base: 12
---
timeline
  2026-Q1 : Project kick-off
  2026-Q2 : Beta release
  2026-Q3 : GA
\end{lstlisting}

\paragraph{Init directive (\texttt{\%\%\{init:\ \ldots\}\%\%}).}

For compatibility with tools that strip frontmatter, the same surface is available via
Mermaid's existing init directive hook:

\begin{lstlisting}[language=yaml,caption={Theme selection via Mermaid init directive}]
%%{init: {"theme": "executive", "layout": "vertical-spine", "density": "compact"}}%%
timeline
  2026-Q1 : Project kick-off
\end{lstlisting}

\paragraph{Surface contract.}
Mermaid is the floor: any standard Mermaid document is valid input, and the compiler
respects Mermaid's built-in theme names (\texttt{default}, \texttt{dark}, \texttt{neutral},
\texttt{forest}, \texttt{base}) as aliases to our nearest equivalent. Our named themes and
token overrides are the ceiling---available but never required. A document with no theming
directives renders under the default theme~\cite{mermaiddocs2024}.

% ──────────────────────────────────────────────────────────────────────────────
\subsection{Determinism Under Theme-Driven Layout}
\label{sec:themes-determinism}

Theme-driven geometry does not compromise determinism. The constraint is:

\begin{quote}
\textbf{Determinism rule (binding):} Themes supply \emph{constants}; engines compute
coordinates \emph{closed-form} from those constants. Same IR + same theme + same layout
argument $\to$ byte-identical output. Themes may change \emph{which} constants are
supplied; they may not introduce randomness, iteration-count sensitivity, or
non-deterministic solver state~\cite{vegalite2017}.
\end{quote}

This is already the case for the existing layout-driving tokens: \texttt{spineSpacing} and
\texttt{nodeWrap} are constants supplied to the layout engine; the engine's coordinate
computation is a closed-form function of those constants. Extending this discipline to the
general contract requires only that every new token be a scalar, enumerated value, or
named step---not a runtime-computed value.

Effect tokens (\texttt{drop-shadow}, \texttt{glow}) degrade gracefully per backend; the
\emph{geometry} under any effect is always identical (see Section~\ref{sec:fidelity-tiers}).

% ──────────────────────────────────────────────────────────────────────────────
\subsection{Fidelity Tiers and Backend Effect Profiles}
\label{sec:fidelity-tiers}

Fidelity tiers define what class of visual richness a theme requests and therefore which
rendering backend is sufficient to honour it fully. A backend that cannot satisfy the
requested tier applies the fallback policies declared in each effect's
\texttt{fallback-policy} field and degrades gracefully; \emph{geometry is never affected}.

\begin{table}[h]
\centering
\small
\begin{tabularx}{\textwidth}{lXXXX}
\toprule
\textbf{Tier} & \textbf{Name} & \textbf{Effect classes} & \textbf{Min.\ backend} & \textbf{Examples} \\
\midrule
0 & Minimal &
  No effects; solid fills, patterns, and opacity only &
  SVG backend &
  Minimal \\[4pt]
1 & Crisp &
  Linear/radial gradients, diagonal hatch, dashed borders, clip paths &
  SVG backend &
  Consulting, Release \\[4pt]
2 & Polished &
  Drop shadows, soft glow; native PPTX effects; SVG safe-filter set (determinism caveat) &
  SVG (caveat) or Raster &
  Executive, Product \\[4pt]
3 & Showcase &
  Bloom, cloud/atmosphere layers, noise textures, gradient meshes, volumetric compositing &
  Raster backend required &
  Showcase, ByteByteGo \\
\bottomrule
\end{tabularx}
\caption{Fidelity tier definitions, effect classes, minimum required backend, and example themes}
\label{tab:fidelity-tiers}
\end{table}

\paragraph{Degradation policy.}
A Tier-3 cloud-layer effect targeting the SVG backend:
\begin{enumerate}
  \item The SVG backend checks capability against \texttt{cloud-layer}: ``not available.''
  \item It reads the effect's \texttt{fallback-policy}: \texttt{embed-raster}.
  \item It invokes the Raster backend for the cloud-layer group only, receives a PNG, and
        embeds it as a \texttt{<image>} element. The rest of the diagram remains clean
        vector.
  \item If the Raster backend is also unavailable, the policy chain falls to
        \texttt{omit}: the cloud layer is silently dropped; no geometry changes.
\end{enumerate}

\paragraph{Geometry invariant under degradation.}
Fallback decisions never alter coordinates, dimensions, or label positions. The Scene
geometry is the invariant shared across all backends and all fidelity levels.

% ──────────────────────────────────────────────────────────────────────────────
\subsection{Theme Inheritance and Extension}
\label{sec:theme-inheritance}

Themes support single-level inheritance via the \texttt{extends} key. A child theme
specifies only the tokens that differ from its parent; all others inherit. This enables
branded variants (\texttt{acme-consulting} extending \texttt{consulting}) with minimal
duplication.

\begin{lstlisting}[language=yaml,caption={Custom theme extending Consulting (general-contract form)}]
theme-id:      "acme-corp"
display-name:  "ACME Corporate"
extends:       "consulting"

palette:
  accent:       "#8B0000"   # ACME brand red replaces navy
  accent-strong: "#5A0000"

typography:
  family:   "Futura"
  fallback: "Century Gothic, sans-serif"
  files:
    regular: "fonts/Futura-Book.woff2"
    bold:    "fonts/Futura-Bold.woff2"
\end{lstlisting}

Inheritance rules:
\begin{enumerate}
  \item Properties not specified in the child are taken from the parent exactly.
  \item \texttt{palette} and nested maps are merged entry-by-entry: a child overrides
        only the entries it declares; unspecified entries inherit from parent.
  \item A font family change requires supplying the corresponding embedded font files.
        A theme that names a font without providing files is malformed.
  \item Inheritance depth is exactly one level. Chains
        (\texttt{A} extends \texttt{B} extends \texttt{C}) are not supported.
  \item The general contract must be fully satisfiable from the resolved (merged) theme.
        A child that leaves a required contract field unresolved is malformed.
\end{enumerate}

% ──────────────────────────────────────────────────────────────────────────────
\subsection{Built-in Named Themes}
\label{sec:builtin-themes}

The compiler ships six built-in named themes that address the primary use cases identified
across David's research and the existing timeline implementation. Each is a complete
realisation of the general contract; each component's implementation produces its own
visual interpretation of the same theme tokens~\cite{thinkcell,mermaid2023}.

\begin{description}
  \item[\texttt{consulting} (Tier 1 — Crisp)] McKinsey/BCG-style. Black, white, single navy
        accent. Heavy typography, generous whitespace, square corners, no decorative chrome.
        \texttt{density: normal}. Optimised for A4/Letter print and A3 poster output.

  \item[\texttt{executive} (Tier 2 — Polished)] Corporate boardroom. Navy/slate on white,
        subtle quarter shading, 4~px corner radius on bars, serif heading font. Status
        communicated by colour $+$ icon. \texttt{density: normal}. 16:9 slide target.

  \item[\texttt{product} (Tier 2 — Polished)] Developer-friendly; information-dense. Per-
        track colour coding. Category colours over status colours; status via small icon.
        \texttt{density: compact}. 1400~px canvas.

  \item[\texttt{release} (Tier 1 — Crisp)] Engineering release calendar. Traffic-light
        palette (green/amber/red). Monospace font. Today marker most prominent element.
        \texttt{density: compact}. Optimised for CI/CD dashboard display.

  \item[\texttt{minimal} (Tier 0 — Minimal)] Maximum whitespace, minimum decoration.
        Single grey fill regardless of status; pattern differentiates status. Print-
        optimised; renders correctly in greyscale. \texttt{density: comfortable}.

  \item[\texttt{showcase} (Tier 3 — Showcase)] Premium keynote aesthetic. Gradient fills,
        drop shadows on nodes, optional atmospheric cloud layer. Full fidelity requires
        Raster backend. \texttt{density: normal}. 1920~px canvas.
\end{description}

\paragraph{Cross-component example.}
Consider the \texttt{consulting} theme applied to three component types simultaneously in a
composition poster:
\begin{itemize}
  \item \textbf{Timeline component:} \texttt{spineSpacing: time}, \texttt{accent} drives
        activity bars, \texttt{corner-radius: 0} produces square bars, \texttt{scale.base}
        = 11~pt for activity labels.
  \item \textbf{Flow component:} \texttt{accent} drives node fills, \texttt{ink} drives
        edge labels, \texttt{corner-radius: 0} produces rectangular nodes, \texttt{scale.sm}
        = 9~pt for edge labels.
  \item \textbf{Sequence component:} \texttt{accent} drives participant header fills,
        \texttt{ink-muted} drives lifeline colour, \texttt{scale.base} = 11~pt for
        message labels.
\end{itemize}
All three diagrams share the same navy accent, the same type scale, the same square shape
language, and the same generous whitespace rhythm---because all three consume the same
\texttt{consulting} general-contract tokens. The poster is visually coherent without any
per-component style co-ordination.

% ──────────────────────────────────────────────────────────────────────────────
\subsection{The Three-Tier Token Architecture}
\label{sec:token-tiers}

Token vocabulary is organized into three tiers following the standard design-token
pattern~\cite{dtcg2024}. The tier boundaries are firm: references flow
\emph{upward only}; the general contract is never contaminated by component-specific
concerns.

\begin{description}
  \item[Tier~1 — Primitives] Raw, un-themed values: colour ramps (HSL ladders, raw hex
        families), typeface family names, the spacing base unit. These are the raw
        material from which Tier-2 tokens are assembled. Components \emph{never} reference
        Tier-1 tokens directly; themes build Tier-2 values from them.

  \item[Tier~2 — Semantic tokens (the general contract)] Role palette, data palette,
        type scale and weights, spacing steps, density, shape language (corner radius,
        node padding, stroke scale, connector style), and effects/fidelity tier.
        \textbf{This is the interface every component implements against.} The full Tier-2
        shape is specified in Section~\ref{sec:themes-contract} and illustrated in the
        listings above.

  \item[Tier~3 — Component tokens] Namespaced \texttt{components.\textit{name}.\textit{token}}
        (e.g.\ \texttt{components.serpentine.turnRadius},
        \texttt{components.sankey.ribbonOpacity},
        \texttt{components.gitgraph.mergeCurveTension}). Each Tier-3 token is
        \emph{derived} from a Tier-2 token by the component's theme-binding as its default.
        A concrete theme \emph{may} override any Tier-3 token as an optional fine-tuning
        step; it is never required to do so.
\end{description}

\begin{quote}
\textbf{Binding invariants (tier contract):}
\begin{itemize}
  \item The general contract (Tier~2) \textbf{NEVER} references a Tier-3 component token.
        The contract is component-agnostic by construction.
  \item Components reference \textbf{upward only}: Tier-3 tokens derive from Tier-2
        defaults; a component may not require a Tier-2 token that does not exist in the
        contract.
  \item A theme \textbf{may} reach downward into a \texttt{components.\textit{name}}
        namespace as an optional override. This is the escape hatch for per-component
        fine-tuning; it is never mandatory.
\end{itemize}
\end{quote}

\subsubsection{Tier-3 Example: serpentine turn radius}

\begin{lstlisting}[language=yaml,caption={Tier-3 token deriving from Tier-2 with optional theme override}]
# ── Tier 2 (general contract, theme-authored) ──────────────────────────
shape:
  corner-radius: 4       # px — general shape language

# ── Tier 3 (component binding, inside serpentine layout engine) ─────────
# Default derivation (computed at bind time, not stored in the theme file):
#   components.serpentine.turnRadius = shape.corner-radius * 6
#   → 4 * 6 = 24 px (sufficient arc to avoid a sharp hairpin)

# ── Optional Tier-3 override in a concrete theme (e.g. "showcase") ──────
components:
  serpentine:
    turnRadius: 40        # override: generous arc for showcase aesthetic
  sankey:
    ribbonOpacity: 0.55   # override: slightly more opaque ribbons
  gitgraph:
    mergeCurveTension: 0.4  # override: softer merge-curve
\end{lstlisting}

The \texttt{components} block is optional in every theme. A theme that omits it entirely
is valid and fully correct: all Tier-3 tokens resolve to their Tier-2 derivation defaults.
Only themes with a strong visual rationale for per-component fine-tuning need to supply it.

% ──────────────────────────────────────────────────────────────────────────────
\subsection{Contract Vocabulary Summary}
\label{sec:contract-summary}

The Tier-2 contract in full, showing all six domains together:

\begin{lstlisting}[language=yaml,caption={Full Tier-2 general contract skeleton (decided)}]
# ── Palette ───────────────────────────────────────────────────────────────────
palette:            # role palette — structural components
  surface: …        # see §12.3 for full role list
  surface-panel: …  # lane / panel / track-header background
  ink: …
  ink-panel: …      # text rendered on panel backgrounds
  accent: …
  border: …
  muted: …
  status-success: { fill: …, stroke: … }
  status-warning: { fill: …, stroke: … }
  status-error:   { fill: …, stroke: … }
  status-info:    { fill: …, stroke: … }
  # workflow states: status-planned, status-active, status-done,
  #                  status-cancelled, status-uncertain
  status-pattern: "solid"  # solid | hatch | dashed-border

data-palette:       # data palette — chart/quantity components
  categorical: [ …, …, …, …, …, … ]   # min 6 entries, ordered
  sequential:  { low: …, mid: …, high: … }
  diverging:   { negative: …, zero: …, positive: … }

# ── Typography ────────────────────────────────────────────────────────────────
typography:
  family: …
  fallback: …
  files: { regular: …, bold: … }
  scale: { xs: 8, sm: 9, base: 11, md: 13, lg: 18, xl: 24 }   # pt
  weights: { regular: 400, medium: 500, semibold: 600, bold: 700 }

# ── Spacing (advisory) ────────────────────────────────────────────────────────
spacing:
  unit: 8   # px — base rhythm unit
  steps: { xxs: 2, xs: 4, sm: 8, md: 16, lg: 24, xl: 32, xxl: 48 }

# ── Density ───────────────────────────────────────────────────────────────────
density: "normal"   # compact | normal | comfortable

# ── Shape language ────────────────────────────────────────────────────────────
shape:
  corner-radius: 0
  node-padding:  8
  stroke-scale:  1.0
  connector-style: "elbow"  # straight | elbow | curved | orthogonal
  marker-shape: "circle"    # circle | diamond | square

# ── Effects / fidelity ────────────────────────────────────────────────────────
effects:
  fidelity: 1   # 0=Minimal | 1=Crisp | 2=Polished | 3=Showcase
  drop-shadow: { enabled: false }
  glow:        { enabled: false }
  motion:      { enabled: false }
\end{lstlisting}

% ──────────────────────────────────────────────────────────────────────────────
\subsection{Contract Adoption Plan}
\label{sec:contract-adoption}

\subsubsection{Proof Set: Validate the Contract First}

Before migrating any production component, the contract is validated by a spike: implement
the \texttt{executive} theme against three components chosen to exercise the entire
contract surface jointly.

\begin{description}
  \item[\texttt{flowchart}] Node-link boxes, edge routing, orientation. Exercises: role
        palette, shape language (corner radius, connector style), type scale, density,
        spacing advisory steps.

  \item[\texttt{sequence}] Text/compartment geometry, activation bars, fragments, lifelines.
        Exercises: role palette, type scale, density (row heights, secondary label
        visibility), shape language (node padding, corner radius).

  \item[\texttt{xychart}] Axes, gridlines, series bars/lines, tick labels. Exercises: the
        \emph{data palette} (categorical sequence + sequential ramp), type scale, spacing
        rhythm. This is the critical third leg: the other two components do not touch the
        data palette, so xychart is required to prove the full contract surface.
\end{description}

The spike succeeds when all three diagrams render coherently under the single
\texttt{executive} theme token set without any component-specific hacks. A passing spike
commits the Tier-2 vocabulary as final~\cite{dtcg2024}.

\subsubsection{Migration Order}

After the proof spike, the 16 existing per-component \texttt{theme.ts} files are migrated
in the following order:

\begin{enumerate}
  \item \textbf{Generalise the timeline \texttt{ResolvedTheme} upward} into the Tier-2
        contract (\texttt{packages/core/src/themes/types.ts}). The timeline is the deepest
        existing layout-driven theme and serves as the reference implementation; all
        subsequent migrations align to the shape it establishes.

  \item \textbf{Node-link family:} \texttt{class}, \texttt{state}, \texttt{ER}, \texttt{C4},
        \texttt{requirement}, \texttt{block}, \texttt{architecture}. These share the same
        fundamental node-edge visual vocabulary and are highest value after the proof set.

  \item \textbf{Remaining charts:} \texttt{pie}, \texttt{quadrant}, \texttt{radar}.
        These consume the data palette and validate categorical and diverging ramp usage
        beyond xychart.

  \item \textbf{Timeline-family adoption of the contract.} The timeline component
        re-consumes the now-general Tier-2 contract, keeping its Tier-3 tokens
        (\texttt{components.serpentine.*}, \texttt{components.roadmap.*}) as derived
        overrides.

  \item \textbf{Specialised components:} \texttt{sankey}, \texttt{gitGraph},
        \texttt{journey}, \texttt{kanban}, \texttt{mindmap}, \texttt{packet}.
        These have the most Tier-3 component-specific tokens and are migrated last
        once the Tier-3 override mechanism is proven by earlier steps.
\end{enumerate}

\begin{quote}
\textbf{Migration invariant:} Each step is independently shippable. No step requires a
downstream step to complete. At every point the compiler is in a valid state: un-migrated
components continue to consume their legacy \texttt{theme.ts}; migrated components consume
the general contract. Both coexist until Step~5 completes.
\end{quote}

% ──────────────────────────────────────────────────────────────────────────────
\subsection{Implementation Status (2026)}
\label{sec:contract-impl-status}

The contract described in this section is fully implemented. The Tier-2
\texttt{ThemeContract} interface and the \texttt{executive} reference theme reside in
\texttt{packages/core/src/theme-contract/}. Every one of the 21 diagram components now
adopts the contract via a per-component binding: for most grammars the binding is
\texttt{grammars/\textit{<component>}/contract-binding.ts}; the timeline uses
\texttt{themes/contract-binding.ts}. Adoption is \textbf{opt-in}: a theme name registered
in \texttt{CONTRACT\_THEMES} resolves via the binding; all legacy per-component named themes
are unchanged and their rendered output is byte-identical to pre-contract builds. The
\texttt{executive} theme has been verified to render all 21 diagram types as one coherent
design system (navy accent, Georgia serif, white surface).

\begin{itemize}
  \item \textbf{Precedence rule:} Legacy component theme names always win. A contract
        theme name must not duplicate a name already claimed by a legacy per-component
        theme; the \texttt{CONTRACT\_THEMES} resolver never overrides a legacy route.
  \item \textbf{Contract location:}
        \texttt{packages/core/src/theme-contract/}
        (\texttt{ThemeContract} Tier-2 interface $+$ \texttt{executive} concrete theme).
  \item \textbf{Binding pattern:}
        \texttt{grammars/\textit{<component>}/contract-binding.ts} (20 grammars)
        and \texttt{themes/contract-binding.ts} (timeline) --- 21 bindings total.
\end{itemize}
